\section{Введение}
\label{sec:Chapter0}

Надёжная локальная навигация является одним из ключевых условий безопасной и
устойчивой работы беспилотных летательных аппаратов (БПЛА) в стеснённых
пространствах и в условиях ухудшенного или отсутствующего приёма сигналов
глобальных навигационных спутниковых систем (GNSS, в том числе ГЛОНАСС). Даже в случае
привязного летательного аппарата, когда полёт механически ограничен тросом,
задача точного и быстрого определения координат остаётся критичной: она
необходима для стабилизации и компенсации внешних возмущений (ветер, натяжение
троса), а также для реализации режимов удержания позиции и сопровождения
объектов. Это особенно важно для БПЛА, несущих на себе телекоммуникационное
оборудование: область покрытия и возможность передачи сигнала напрямую зависят
от положения летательного аппарата~\cite{vishnevsky2013}.

На практике применяются различные подходы к локализации БПЛА: спутниковая
навигация (в том числе дифференциальные и RTK-решения), визуальные и
визуально-инерциальные методы, лидарная одометрия, инфраструктурные системы
захвата движения. Каждый из них имеет свою область применимости и свои
ограничения; подробное сравнение этих подходов применительно к рассматриваемой
задаче приведено в разделе~\ref{sec:Chapter2}.

Ультразвуковая навигация представляет собой привлекательный вариант для задач
ближнего радиуса действия благодаря низкой стоимости, малому энергопотреблению и
доступности компактных излучателей и приёмников. В таких системах положение
объекта оценивается по времени распространения акустического сигнала между
передатчиками и приёмниками с последующим пересчётом временных измерений в
расстояния на основе величины скорости звука. В зависимости от архитектуры могут
применяться схемы измерения времени пролёта (ToF~--- от англ.\ \textit{Time of
Flight}), разности времён прихода (TDoA~--- от англ.\ \textit{Time Difference of
Arrival}) или их комбинации с геометрическими методами трилатерации и
мультилатерации. В отличие от чисто бортовых методов измерения, ультразвуковые
системы позволяют опираться на внешние опорные узлы (якоря) и могут сохранять работоспособность
даже в тех условиях, где оптика даёт сбои. Вместе с тем ультразвуковое
позиционирование имеет ограничения: сравнительно малую дальность и
чувствительность к отражениям сигнала от поверхностей.

В настоящей работе рассматривается задача локальной навигации привязного БПЛА
относительно платформы (в том числе потенциально подвижной) с использованием
ультразвуковых измерений в качестве основного источника координатной информации.
Ставится цель разработать систему, пригодную для практической реализации,
обосновать её основные параметры математической моделью и подтвердить
работоспособность экспериментально.
